\section{Light --- Frames}
%%% generated by the blueprint pipeline from TCSlib.Programs.Light.Frames
\subsection*{Light --- Frames: Introduction}

Declarations of the Lean module \texttt{TCSlib.Programs.Light.Frames}.

\subsection*{Light --- Frames: Results}

\begin{definition}[Update one local variable]
\lean{Light.setLocal}
\label{Light.setLocal}
\leanok
Given a list of integer local-variable values, a natural-number index $x$, and an
integer $z$, define the updated list by replacing the value at index $x$ with $z$. If
the list is too short, extend it with zeros as needed to reach index $x$, then place $z$
there.
\end{definition}

\begin{theorem}[Updating a local frame entry]
\lean{Light.frame_setLocal}
\label{Light.frame_setLocal}
\leanok
\uses{Light.frame, Light.setLocal}
\difficulty{4}
For a finite list $l$ of integers, an index $x\in\mathbb N$, and an integer $z$, let the
updated list be obtained by replacing its entry at $x$ by $z$, padding with zeros if
necessary. At every index $y$, the frame of the updated list equals $z$ if $y=x$, and
otherwise equals the frame of $l$.
\end{theorem}

\begin{theorem}[Updating a list frame at one index]
\lean{Light.update_frame_setLocal}
\label{Light.update_frame_setLocal}
\leanok
\uses{Light.frame, Light.frame_setLocal, Light.setLocal}
\difficulty{3}
For every integer list $l$, natural-number index $x$, and integer $z$, updating the
function represented by the zero-padded list $l$ at index $x$ to have value $z$ gives
the function represented by the list obtained by setting entry $x$ of $l$ to $z$,
extending the list with zeros as needed.
\end{theorem}

\begin{theorem}[Appending zeros preserves the frame]
\lean{Light.frame_append_zeros}
\label{Light.frame_append_zeros}
\leanok
\uses{Light.frame}
\difficulty{4}
For every finite list $l$ of integers and natural number $n$, the function that reads
the list at an index and returns zero when the index is out of range is unchanged when
$n$ zeros are appended to $l$.
\end{theorem}

\begin{definition}[Expression gives a value]
\lean{Light.Expr.Gives}
\label{Light.Expr.Gives}
\leanok
\uses{Light.Expr, Light.Expr.Safe, Light.Expr.val, Light.Limits, Light.State}
For resource limits $\text{lim}$, a state $\sigma$, an expression $e$, and an integer
$z$, say that $e$ gives $z$ when $e$ is safe under $\text{lim}$ in state $\sigma$ and
evaluates there to $z$.
\end{definition}

\begin{theorem}[Assignment with a postcondition]
\lean{Light.Ends.setVal}
\label{Light.Ends.setVal}
\leanok
\uses{Light.Ends, Light.Ends.set, Light.Expr, Light.Expr.Gives, Light.Expr.cost}
\difficulty{1}
Let a limit record, program, execution depth, local-memory function $\mu$, budget $T$,
postcondition $Q$, local function $\mathrm{loc}$, index $x$, expression $e$, and integer
$z$ be given. If $Q$ holds of the state with local function obtained by updating
$\mathrm{loc}$ at $x$ to $z$ and memory $\mu$, and $e$ is safe in the state
$(\mathrm{loc},\mu)$ and evaluates there to $z$, with $\mathrm{cost}(e)+1\le T$, then
assigning $e$ to position $x$ can execute from $(\mathrm{loc},\mu)$ with cost at most
$T$ and yield a state satisfying $Q$.
\end{theorem}

\begin{theorem}[Assignment followed by a program]
\lean{Light.Ends.setValThen}
\label{Light.Ends.setValThen}
\leanok
\uses{Light.Ends, Light.Ends.next, Light.Ends.setVal, Light.Expr, Light.Expr.Gives, Light.Expr.cost, Light.Stmt}
\difficulty{1}
Let a limit record, program, execution depth, memory function $\mu$, budget $T$,
postcondition $Q$, local function $\mathrm{loc}$, index $x$, expression $e$, statement
$s$, and integer $z$ be given. If $s$ can execute from the state obtained by updating
local position $x$ to $z$ in $(\mathrm{loc},\mu)$, with cost at most
$T-(\mathrm{cost}(e)+1)$, and its final state satisfies $Q$, and if $e$ is safe in
$(\mathrm{loc},\mu)$ and evaluates there to $z$, with $\mathrm{cost}(e)+1\le T$, then
assigning $e$ to position $x$ and then executing $s$ can execute from
$(\mathrm{loc},\mu)$ with cost at most $T$ and final state satisfying $Q$.
\end{theorem}

\begin{theorem}[Assignment to a prescribed local value]
\lean{Light.Ends.setTo}
\label{Light.Ends.setTo}
\leanok
\uses{Light.Ends, Light.Ends.set, Light.Expr, Light.Expr.Gives, Light.Expr.cost, Light.frame, Light.setLocal, Light.update_frame_setLocal}
\difficulty{3}
Let $l$ be a finite list of integers, $\mu:\mathbb{N}\to\mathbb{Z}$ a memory function,
$x$ a natural-number local index, $z$ an integer, $e$ an expression, $T$ a time bound,
and $Q$ a predicate on states. Suppose $Q$ holds in the state with memory $\mu$ and
local frame obtained by setting entry $x$ of $l$ to $z$. Suppose $e$ is safe in the state
with local frame $l$ and memory $\mu$, evaluates there to $z$, and has cost satisfying
$\mathrm{cost}(e)+1\le T$. Then assigning $e$ to local variable $x$ from that state
ends within time $T$ in a state satisfying $Q$.
\end{theorem}

\begin{theorem}[Set followed by a program]
\lean{Light.Ends.setToThen}
\label{Light.Ends.setToThen}
\leanok
\uses{Light.Ends, Light.Ends.next, Light.Ends.setTo, Light.Expr, Light.Expr.Gives, Light.Expr.cost, Light.Stmt, Light.frame, Light.setLocal}
\difficulty{1}
Let a limit record, program, execution depth, list of integer values $l$, local-memory
function $\mu$, budget $T$, postcondition $Q$, index $x$, expression $e$, statement $s$,
and integer $z$ be given. If $s$ can execute from the state whose local values are
obtained by setting position $x$ of the frame of $l$ to $z$, using cost at most
$T-(\mathrm{cost}(e)+1)$, and its final state satisfies $Q$, and if $e$ is safe in the
state with local values given by the frame of $l$ and memory $\mu$ and evaluates there
to $z$, with $\mathrm{cost}(e)+1\le T$, then the sequence that first assigns $e$ to
position $x$ and then executes $s$ can execute from that initial state with cost at most
$T$ and final state satisfying $Q$.
\end{theorem}

\begin{theorem}[Store to a specified address and value]
\lean{Light.Ends.storeTo}
\label{Light.Ends.storeTo}
\leanok
\uses{Light.Ends, Light.Ends.store, Light.Expr, Light.Expr.Gives, Light.Expr.cost, Light.frame}
\difficulty{4}
Let a program, depth bound, list $l$ of integers, memory map $\mu:\mathbb N\to\mathbb
Z$, time bound $T\in\mathbb N$, postcondition $Q$, expressions $a,e$, address
$b\in\mathbb N$, and value $z\in\mathbb Z$ be given. If $Q$ holds in the state with
locations given by the frame of $l$ and memory updated at $b$ to $z$, if $a$ and $e$ are
safe in the state with those locations and original memory and evaluate respectively to
$b$ and $z$, and if $b$ is a valid address, then storing $e$ at address $a$ ends in a
state satisfying $Q$ within time $T$, provided the store cost
$\operatorname{cost}(a)+\operatorname{cost}(e)+1$ is at most $T$.
\end{theorem}

\begin{theorem}[Store followed by a bounded computation]
\lean{Light.Ends.storeToThen}
\label{Light.Ends.storeToThen}
\leanok
\uses{Light.Ends, Light.Ends.next, Light.Ends.storeTo, Light.Expr, Light.Expr.Gives, Light.Expr.cost, Light.Stmt, Light.frame}
\difficulty{1}
Let a program, depth bound, list $l$ of integers, memory map $\mu:\mathbb N\to\mathbb
Z$, time bound $T\in\mathbb N$, postcondition $Q$, expressions $a,e$, statement $s$,
address $b\in\mathbb N$, and value $z\in\mathbb Z$ be given. Suppose that $s$ can
execute from the state with locations given by the frame of $l$ and memory updated at
$b$ to $z$, within time $T-(\operatorname{cost}(a)+\operatorname{cost}(e)+1)$, and end
in a state satisfying $Q$. Suppose also that $a,e$ are safe in the original state and
evaluate respectively to $b,z$, and that $b$ is less than the space bound. Then storing
$e$ at address $a$ and executing $s$ from the original state ends in a state satisfying
$Q$ within time $T$, provided the store cost is at most $T$.
\end{theorem}

\begin{theorem}[For-loop over a framed local list]
\lean{Light.Ends.forFrame}
\label{Light.Ends.forFrame}
\leanok
\uses{Light.Ends, Light.Ends.forMem, Light.Expr, Light.Expr.Gives, Light.Expr.cost, Light.Stmt, Light.Stmt.blockCost, Light.Stmt.for, Light.frame, Light.setLocal, Light.update_frame_setLocal}
\difficulty{4}
Let $\mathrm{lim}$ be execution limits, $P$ a program, $l$ a list of integers, $\mu$ an
initial memory, $i$ a local index, $hi$ an upper-bound expression, $body$ a loop body, and
$Q$ a postcondition. Let $n,T$ be natural numbers and $I(j,\mu\prime)$ an invariant on
a loop index and memory. Assume $I(0,\mu)$. For every $j<n$ and memory $\mu\prime$
satisfying $I(j,\mu\prime)$, suppose the body executes from the state with local frame
$l$ updated at $i$ to $j$ and memory $\mu\prime$, within its block-cost bound, and
ends with that same local frame and a memory satisfying $I(j+1,\cdot)$. Suppose also
that $I(n,\mu\prime)$ implies $Q$ for the state with index $i$ set to $n$, and that
whenever $j\le n$ and $I(j,\mu\prime)$ holds, $hi$ is safe and evaluates to $n$ in
the state with index $i$ set to $j$. If $n\le\mathrm{lim.word}$ and
$n(\mathrm{cost}(hi)+\mathrm{blockCost}(body)+7)+\mathrm{cost}(hi)+5\le T$, then the
for-loop at index $i$ terminates within time $T$ from the state with local frame $l$
and memory $\mu$, ending in a state satisfying $Q$.
\end{theorem}

\begin{theorem}[Bounded execution of a shaped loop]
\lean{Light.Ends.forShape}
\label{Light.Ends.forShape}
\leanok
\uses{Light.Ends, Light.Ends.for, Light.Ends.mono, Light.Expr, Light.Expr.Gives, Light.Expr.cost, Light.State, Light.Stmt, Light.Stmt.for}
\difficulty{5}
Fix a program environment, resource limits, and execution depth. Let $S(j,s,\mu)$ be a
state for a counter value $j$, shape parameter $s$, and memory $\mu$, and let $I(j,\mu)$
be an invariant. Suppose $I(0,\mu)$ holds, and for each $j<n$, each $s$, and each $\mu'$
satisfying $I(j,\mu')$, the body can execute from $S(j,s,\mu')$ with cost at most $b$ to
a state of the form $S(j,s',\mu'')$ with $I(j+1,\mu'')$. Suppose $Q(S(n,s,\mu'))$ holds
for every $s,\mu'$ satisfying $I(n,\mu')$. Also suppose the state obtained by setting
local $i$ to zero in the initial local store $\ell$ and using memory $\mu$ equals
$S(0,s_0,\mu)$; for every $j\le n$ and every $s,\mu'$ satisfying $I(j,\mu')$, the bound
expression is safe in $S(j,s,\mu')$ and evaluates there to $n$; the value of local $i$
in each $S(j,s,\mu')$ is $j$; and setting local $i$ to $j+1$ in $S(j,s,\mu')$ gives
$S(j+1,s,\mu')$. If $n\le\mathrm{word}$ and
$n(\mathrm{cost}(h)+b+7)+\mathrm{cost}(h)+5\le T$, then the counted loop starting from
the state with locals $\ell$ and memory $\mu$ can finish with cost at most $T$ in a
state satisfying $Q$.
\end{theorem}

\begin{definition}[Update a list of local variables]
\lean{Light.setLocals}
\label{Light.setLocals}
\leanok
\uses{Light.setLocal}
Given a list $l$ of integer local-variable values and a list of index-value pairs,
define the resulting list by applying each update in order: at each pair $(x,z)$,
replace the value at index $x$ by $z$, extending the list with zeros when necessary.
\end{definition}

\begin{definition}[Updating local values]
\lean{Light.updateLocals}
\label{Light.updateLocals}
\leanok
Given a function assigning an integer to each natural-number index and a list of
index–integer pairs, define the updated function by applying the listed assignments in
order, so later assignments to an index determine its final value.
\end{definition}

\begin{theorem}[Updating a framed list of local values]
\lean{Light.updateLocals_frame}
\label{Light.updateLocals_frame}
\leanok
\uses{Light.frame, Light.setLocals, Light.updateLocals, Light.update_frame_setLocal}
\difficulty{4}
For every finite list $ps$ of pairs of a natural-number index and an integer value, and
every integer list $l$, updating the function represented by the zero-padded list $l$ at
each index-value pair in $ps$ gives the function represented by the list obtained by
applying those updates successively to $l$.
\end{theorem}

\begin{definition}[Refreshing local values]
\lean{Light.refreshLocals}
\label{Light.refreshLocals}
\leanok
\uses{Light.setLocal}
Given an integer list, a function assigning an integer to each natural-number index, and
a list of indices, define the refreshed list by setting each listed position to the
function’s value at that index, in order. Setting a position beyond the current list
length extends the list with zeros as needed.
\end{definition}

\begin{theorem}[Refreshing selected local values]
\lean{Light.frame_refreshLocals}
\label{Light.frame_refreshLocals}
\leanok
\uses{Light.frame, Light.frame_setLocal, Light.refreshLocals}
\difficulty{4}
Let $loc:\mathbb N\to\mathbb Z$, let $y\in\mathbb N$, and let $xs$ be a finite list of
natural numbers. Refresh a finite list of integer values by setting, in sequence, the
entry at each index in $xs$ to its value under $loc$, extending the list with zeros as
needed. The resulting list read at index $y$ equals $loc(y)$ if $y$ occurs in $xs$, and
otherwise equals the original list read at $y$.
\end{theorem}

\begin{definition}[Equality of frames]
\lean{Light.SameFrame}
\label{Light.SameFrame}
\leanok
\uses{Light.frame}
Two integer lists have the same frame if their frames are equal.
\end{definition}

\begin{theorem}[Equality of empty frames]
\lean{Light.sameFrame_nil}
\label{Light.sameFrame_nil}
\leanok
\uses{Light.SameFrame}
\difficulty{1}
The zero-padded frames of two empty lists are equal.
\end{theorem}

\begin{theorem}[Equality of zero-padded list frames]
\lean{Light.sameFrame_cons_cons}
\label{Light.sameFrame_cons_cons}
\leanok
\uses{Light.SameFrame}
\difficulty{4}
For integers $a,b$ and lists $l,m$ of integers, the zero-padded frames of $a::l$ and
$b::m$ are equal if and only if $a=b$ and the zero-padded frames of $l$ and $m$ are
equal.
\end{theorem}

\begin{theorem}[Equality of a frame with the empty frame]
\lean{Light.sameFrame_cons_nil}
\label{Light.sameFrame_cons_nil}
\leanok
\uses{Light.SameFrame}
\difficulty{4}
For an integer $a$ and a list $l$ of integers, the zero-padded frames of $a::l$ and the
empty list are equal if and only if $a=0$ and the zero-padded frames of $l$ and the
empty list are equal.
\end{theorem}

\begin{theorem}[Equality of the empty frame with a list frame]
\lean{Light.sameFrame_nil_cons}
\label{Light.sameFrame_nil_cons}
\leanok
\uses{Light.SameFrame}
\difficulty{4}
For an integer $b$ and a list $m$ of integers, the zero-padded frames of the empty list
and $b::m$ are equal if and only if $b=0$ and the zero-padded frames of the empty list
and $m$ are equal.
\end{theorem}

\begin{definition}[Agreement outside a list of indices]
\lean{Light.LocalsBut}
\label{Light.LocalsBut}
\leanok
\uses{Light.frame}
An integer-valued function agrees with the frame of an integer list outside a specified
list of natural-number indices if, for every index not in that list, the function’s
value equals the frame’s value there.
\end{definition}

\begin{theorem}[Frame after refreshing local variables]
\lean{Light.LocalsBut.eq_frame}
\label{Light.LocalsBut.eq_frame}
\leanok
\uses{Light.LocalsBut, Light.frame, Light.frame_refreshLocals, Light.refreshLocals}
\difficulty{3}
For a list $l$ of integers, define its frame at index $i$ to be its entry at $i$, or $0$
if there is no such entry. Let $\mathrm{loc}:\mathbb N\to\mathbb Z$ agree with the frame
of $l$ at every index outside a list $xs$. Then $\mathrm{loc}$ is exactly the frame of
the list obtained by refreshing the entries of $l$ at the indices in $xs$ using the
corresponding values of $\mathrm{loc}$.
\end{theorem}

\begin{theorem}[Updating a local variable]
\lean{Light.LocalsBut.update}
\label{Light.LocalsBut.update}
\leanok
\uses{Light.LocalsBut, Light.frame_setLocal, Light.setLocal}
\difficulty{3}
For a list $l$ of integers, let its frame give the list entry at each index and $0$
beyond the list. Let $\mathrm{loc}:\mathbb N\to\mathbb Z$ agree with the frame of $l$ at
every index outside a list $xs$. For any index $x$ and integer $z$, the function
obtained by setting $\mathrm{loc}(x)=z$ agrees with the frame of the list obtained by
setting its entry at $x$ to $z$, at every index outside $xs$.
\end{theorem}

\begin{theorem}[Agreement outside refreshed indices]
\lean{Light.LocalsBut.of_eq}
\label{Light.LocalsBut.of_eq}
\leanok
\uses{Light.LocalsBut, Light.SameFrame, Light.frame, Light.frame_refreshLocals, Light.refreshLocals}
\difficulty{4}
For a list $l$ of integers, let its frame give the list entry at each index and give $0$
beyond the list. Let $xs$ be a list of indices and $l'$ another list. If refreshing $l$
at the indices in $xs$ using the frame of $l'$ produces a list whose frame equals the
frame of $l'$, then the frame of $l'$ agrees with the frame of $l$ at every index
outside $xs$.
\end{theorem}

\begin{theorem}[Execution from refreshed local variables]
\lean{Light.Ends.asFrame}
\label{Light.Ends.asFrame}
\leanok
\uses{Light.Ends, Light.LocalsBut, Light.LocalsBut.eq_frame, Light.frame, Light.refreshLocals}
\difficulty{1}
Let $xs$ be a list of variable indices, $l$ a list of integers, and $loc$ and $\mu$
integer-valued functions on the natural numbers. Suppose $loc(y)$ equals the value of
the list frame $l$ at index $y$ for every $y$ not in $xs$. If a statement $s$ has an
execution from the state whose locals are obtained by refreshing the entries of $l$ at
the indices in $xs$ with their values from $loc$, and whose memory is $\mu$, within cost
$T$ and with postcondition $Q$, then it has such an execution from the state with locals
$loc$ and memory $\mu$.
\end{theorem}

\begin{theorem}[Execution with refreshed local variables]
\lean{Light.Ends.pieceTo}
\label{Light.Ends.pieceTo}
\leanok
\uses{Light.Ends, Light.Ends.keeping, Light.Ends.mono, Light.LocalsBut, Light.LocalsBut.eq_frame, Light.Stmt.assigns, Light.frame, Light.refreshLocals}
\difficulty{4}
Let $xs$ be a list of local indices, let $s$ be a statement, and let $l$ be a list of
integers. If $s$ has an execution from the state with frame $l$ and memory $\mu$, of
cost at most $T_1$, ending in a state satisfying $R$, and every such final state with
locals $\mathrm{loc}'$ and memory $\mu'$ gives $Q$ in the state with frame obtained by
refreshing the locals in $l$ at indices $xs$ from $\mathrm{loc}'$ and memory $\mu'$,
then $s$ has an execution from the original state with cost at most $T$ ending in a
state satisfying $Q$, provided every local assigned by $s$ occurs in $xs$ and $T_1\le
T$.
\end{theorem}

\begin{theorem}[Sequential composition with refreshed locals]
\lean{Light.Ends.pieceToThen}
\label{Light.Ends.pieceToThen}
\leanok
\uses{Light.Ends, Light.Ends.next, Light.Ends.pieceTo, Light.Stmt.assigns, Light.frame, Light.refreshLocals}
\difficulty{3}
Let $\mathrm{lim}$ be resource limits, $P$ a program, $d$ a call depth, and $s_1,s_2$
statements. Let $l$ be a list of integers, $\mu:\mathbb{N}\to\mathbb{Z}$ a memory, and
$T,T_1$ time bounds. Suppose $s_1$ executes from the state with local frame $l$ and memory
$\mu$ within time $T_1$, ending in a state satisfying $R$. For every state with local
function $\mathrm{loc}\prime$ and memory $\mu\prime$ satisfying $R$, suppose $s_2$ executes
from the state with local frame $l$ refreshed at the indices in $xs$ from $\mathrm{loc}\prime$
and memory $\mu\prime$, within time $T-T_1$, ending in a state satisfying $Q$. If the
indices assigned by $s_1$ are all in $xs$ and $T_1\le T$, then $s_1;s_2$ executes from
the initial state within time $T$ and ends in a state satisfying $Q$.
\end{theorem}

\begin{theorem}[Bounded execution of a counted loop]
\lean{Light.Ends.forScratch}
\label{Light.Ends.forScratch}
\leanok
\uses{Light.Ends, Light.Ends.for, Light.Ends.keeping, Light.Ends.mono, Light.Expr, Light.Expr.Gives, Light.Expr.cost, Light.LocalsBut, Light.State, Light.Stmt, Light.Stmt.assigns, Light.Stmt.for, Light.frame, Light.frame_refreshLocals, Light.frame_setLocal, Light.refreshLocals, Light.setLocal}
\difficulty{6}
Fix a program environment, resource limits, and execution depth, and let $i$ be a local
variable, $h$ an expression, and $B$ a loop body. Let $X$ be a list of local indices,
$I(j,\mu)$ an invariant on a counter value $j$ and memory $\mu$, and $Q$ a postcondition
on states. Suppose $I(0,\mu)$ holds; for every $j<n$, every local store $\ell'$, and
every memory $\mu'$ satisfying $I(j,\mu')$, executing $B$ from the state whose locals
are the frame of the list obtained by refreshing $\ell'$ along $X$ and then setting
local $i$ to $j$, with memory $\mu'$, can finish with cost at most $b$ and memory
satisfying $I(j+1,\cdot)$. Suppose also that whenever $I(n,\mu')$ holds, $Q$ holds of
that same framed state with local $i$ set to $n$; that for every $j\le n$ satisfying the
invariant, $h$ is safe in that state and has value $n$; that every local assigned by $B$
belongs to $X$; that $i\notin X$; that $n\le\mathrm{word}$; and that
$n(\mathrm{cost}(h)+b+7)+\mathrm{cost}(h)+5\le T$. Then the counted loop starting from
the state with locals given by the frame of $\ell$ and memory $\mu$ can finish with cost
at most $T$ in a state satisfying $Q$.
\end{theorem}
